\chapter{[WIS] Wissenschaft und Forschung}
\label{chp:WIS}

\section{Von der Studie zur Lehrmeinung}

\subsubsection{Fachgesellschaften}

In der Welt der Wissenschaft enstehen und entstanden sog. Fachgesellschaften,
die sich bestimmten Fachgebieten, Fragestellungen, Erkrankungen, etc. widmen.

\gLayout{0}{Medizinische Fachgesellschaften}
{
\index{topic:Fachgesellschaft}
Die Tätigkeit von medizinische Fachgesellschaften umfasst oft die Forschung,
Fortbildung und die Herausgabe von
\gEs{Empfehlungen} und Leittlinien zur Diagnostik, Therapie und Vorbeugung
u.ä. Diese Empfehlungen sind nicht "`gesetzlich bindend"', aber geben i.d.R.
den \gEs{Stand der Wissenschaft} wieder. Ausserdem müssen diese, oft allgemein
formulierten, Empfehlungen an die örtlichen Gegebenheiten angepasst
werden. Verschiedene Fachgesellschaften können -- auch zu den gleichen Themen --
unterscheidliche Meinungen haben und unterschiedliche Empfehlungen aussprechen.
}{
\begin{itemize}
   \item Forschung, Fortbildung. 
   \item Herausgabe von Empfehlungen für Diagnostik, Therapie und
   Vorbeugung (u.a.).
   \item Nicht bindend, aber "`Stand der Wissenschaft"'.
\end{itemize}
}{}{}{}


\gLayout{0}{ERC}
{
\index{topic:ERC}
\index{topic:Fachgesellschaft!ERC}
Eine wichtige Fachgesellschaft ist das \gT{ERC} (\gE{European Resuscitation
Council}). Bekannt ist das \gAbk{ERC} u.a. durch seine Leitlinien zur
Wiederbelebung. Alle 5 Jahre werden diese "`Guidelines\footnote{engl.
\emph{Guideline}: Richtlinie, Empfehlung}"' aktualisiert. Jeder neuen Auflage
der Richtlinien geht ein aufwendiger Forschungs- und Diskussionsprozess voraus: Es
werden hunderte Studien mit vielen verschiedenen Fragestellungen an vielen
Forschungszentren auf der ganzen Welt durchgeführt, die Ergebnisse ausgewertet und anschließend diskutiert. 

Die derzeit aktuelle Version
der ERC-Guidelines kam im Jahr 2005 heraus("`\gE{ERC 2005}"'), die nächste
Aktualisierung ist mit Ende des Jahre 2010 zu erwarten.
}{
\begin{itemize}
   \item Leitlinien zur Reanimation.
   \item Aktualisierung alle 5 Jahre.
   \item Letzte 2005, nächste Ende 2010.
\end{itemize}

}{}{}{}


\section{Statisktik}


\gCaptionof{table}{Statistische Tests mit lustigen Namen -- Eine Übersicht.}